\documentclass[11pt,a4paper]{article}
\usepackage[utf8]{inputenc}
\usepackage[T1]{fontenc}
\usepackage[french]{babel}
\usepackage{amsmath}
\usepackage{amssymb}
\usepackage{graphicx}
\usepackage{xcolor}
\usepackage{enumitem}
\usepackage{geometry}
\usepackage{listings} % Pour afficher du code Python proprement
\geometry{left=1.5cm,right=1.5cm,top=1.5cm,bottom=1.5cm}

% Configuration de listings pour Python
\lstset{
    language=Python,
    basicstyle=\ttfamily\small,
    keywordstyle=\color{blue},
    commentstyle=\color{gray},
    stringstyle=\color{red},
    showstringspaces=false,
    breaklines=true,
    frame=single,
    numbers=left,
    numberstyle=\tiny\color{gray},
    xleftmargin=15pt
}

% mettre dans le préambule, après \usepackage{xcolor}
\usepackage{sectsty} % gère facilement la couleur des titres
\setlength{\parindent}{0pt}

\sectionfont{\color{red}\Large\bfseries\underline}     % pour \section*
\subsectionfont{\color{red}\large\bfseries\hspace{1cm}}

\title{Fiche Résumé : Algorithmique et Programmation en Python}
\author{}
\date{}

\begin{document}

\begin{center}
    {\LARGE \textbf{Python -- Algorithmique - Fiche résumée}}\\[4mm]
    {\large -- 1\up{ère} STMG --}
\end{center}

\vspace{20pt}

\section*{I/ Les bases du langage Python}
\subsection*{1) Affectation}
\begin{itemize}[leftmargin=*]
    \item L'\textbf{affectation} permet de stocker une valeur dans une variable.
    \item Syntaxe : \lstinline|variable = valeur|
    \item Exemple : \lstinline|x = 5| ou \lstinline|nom = "Alice"|
\end{itemize}

\subsection*{2) Conditions}
\begin{itemize}[leftmargin=*]
    \item Les conditions permettent d'exécuter du code en fonction d'une condition.
    \item Syntaxe :
    \begin{lstlisting}
if condition:
    # Code execute si la condition est vraie
elif autre_condition:
    # Code execute si la premiere condition est fausse
else:
    # Code execute si aucune condition n'est vraie
    \end{lstlisting}
    \item Opérateurs de comparaison : \lstinline|==|, \lstinline|!=|, \lstinline|>|, \lstinline|<|, \lstinline|>=|, \lstinline|<=|
    \item Opérateurs logiques : \lstinline|and|, \lstinline|or|, \lstinline|not|
\end{itemize}

\section*{II. Les boucles}
\subsection*{1) Boucle \texttt{while}}
\begin{itemize}[leftmargin=*]
    \item La boucle \texttt{while} exécute un bloc de code \textbf{tant qu'une condition est vraie}.
    \item Syntaxe :
    \begin{lstlisting}
while condition:
    # Code execute tant que la condition est vraie
    \end{lstlisting}
\end{itemize}

\subsection*{2) Boucle \texttt{for}}
\begin{itemize}[leftmargin=*]
    \item La boucle \texttt{for} permet d'itérer sur une séquence (liste, plage de nombres, etc.).
    \item Syntaxe :
    \begin{lstlisting}
for variable in sequence:
    # Code execute pour chaque element
    \end{lstlisting}
    \item Exemple : \lstinline|for i in range(5):| (itère de 0 à 4).
    \item \lstinline|for i in range(2, 13, 2)| : (itère de 2 à 12 de 2 en 2).
\end{itemize}


\section*{III/ Les fonctions}
\subsection*{1) Définition et appel}
\begin{itemize}[leftmargin=*]
    \item Une \textbf{fonction} est un bloc de code réutilisable.
    \item Syntaxe :
    \begin{lstlisting}
def nom_fonction(arguments):
    # Code de la fonction
    return valeur  # Optionnel
    \end{lstlisting}
    \item Exemple :
    \begin{lstlisting}
def carre(x):
    return x ** 2
    \end{lstlisting}
    \item Appel : \lstinline|carre(5)| retourne 25.
\end{itemize}

\subsection*{2) Arguments}
\begin{itemize}[leftmargin=*]
    \item Les \textbf{arguments} sont les valeurs passées à une fonction.
    \item Exemple : Dans \lstinline|def somme(a, b):|, \lstinline|a| et \lstinline|b| sont des arguments.
\end{itemize}


\section*{IV/ Génération de nombres aléatoires}
\subsection*{1) Nombres entiers aléatoires}
\begin{itemize}[leftmargin=*]
    \item \lstinline|randint(a, b)| : Génère un entier aléatoire entre \texttt{a} et \texttt{b} (inclus).
    \item Exemple : \lstinline|randint(1, 6)| simule un lancer de dé.
\end{itemize}

\subsection*{2) Nombres entre 0 et 1}
\begin{itemize}[leftmargin=*]
    \item \lstinline|random()| : Génère un nombre flottant aléatoire entre 0 et 1.
\end{itemize}

\section*{V/ Les listes}
\subsection*{1) Création et manipulation}
\begin{itemize}[leftmargin=*]
    \item Une \textbf{liste} est une collection ordonnée d'éléments.
    \item Syntaxe : \lstinline|ma_liste = [1, 2, 3, "quatre"]|
    \item Accès aux éléments : \lstinline|ma_liste[0]| (premier élément), 
    \lstinline|ma_liste[1]| (deuxième élément), etc.
\end{itemize}

\subsection*{2) Ajout et suppression}
\begin{itemize}[leftmargin=*]
    \item \lstinline|ma_liste.append(x)| : Ajoute \texttt{x} à la fin de la liste.
    \item \lstinline|ma_liste.remove(x)| : Supprime la première occurrence de \texttt{x}.
    \item \lstinline|del ma_liste[i]| : Supprime l'élément à l'indice \texttt{i}.
\end{itemize}

\subsection*{3) Itérer sur une liste}
\begin{itemize}[leftmargin=*]
    \item Utiliser une boucle \texttt{for} :
    \begin{lstlisting}
for element in ma_liste:
    print(element)
    \end{lstlisting}
\end{itemize}


\end{document}